% !PRINTED 226
\weberpage{226}
Now by \S.\,68, (6),
\[
S f_0 = n a_0, \quad S f_1 = (n-1) a_1, \ldots S f_{n-1} = a_{n-1},
\]
so that one obtains
\[
\begin{aligned}
S(y z) = {} & n a_0 \sum_{0,\,n-1}^{s} E_{0,s} \tau_{n-s-1} \\
& + (n-1) a_1 \sum_{0,\,n-1}^{s} E_{1,s} \tau_{n-s-1} \\
& \cdots\cdots\cdots\cdots \\
& + a_{n-1} \sum_{0,\,n-1}^{s} E_{n-1,s} \tau_{n-s-1}.
\end{aligned}
\]
In this, now, the coefficient of $\tau_{n-1}$ is
\[
n a_0 E_{0,0} + (n-1) a_1 E_{1,0} + \cdots + a_{n-1} E_{n-1,0},
\]
hence by (3)
\[
= a_0 S(y).
\]
Now since $S(z)$ begins with the term $n a_0 \tau_{n-1}$, in the
difference
\weq{5}{S(y z) - \frac{1}{n} S(y) S(z)}
no term will occur that is multiplied by $\tau_{n-1}$. But since this
expression is moreover symmetric with respect to $t$ and $\tau$,
it also does not contain $t_{n-1}$, and in forming it we may
% [sic] The print omits “=0” after the first term; the immediately following paragraph confirms that both t_{n-1} and tau_{n-1} are set equal to zero.
simply assume $t_{n-1}$ and $\tau_{n-1} = 0$.

We now set
\[
S(y) = n a_0 t_{n-1} + \varphi(t),
\]
in which
\weq{6}{\varphi(t) = (n-1) a_1 t_{n-2} + (n-2) a_2 t_{n-3} + \cdots + a_{n-1} t_0.}
Then we have
\weq{7}{
\begin{aligned}
S(y z) - \frac{1}{n} S(y) S(z) = {} & n a_0 \sum_{1,\,n-1}^{s} E_{0,s} \tau_{n-s-1} \\
& + (n-1) a_1 \sum_{1,\,n-1}^{s} E_{1,s} \tau_{n-s-1} \\
& \cdots\cdots\cdots\cdots \\
& + a_{n-1} \sum_{1,\,n-1}^{s} E_{n-1,s} \tau_{n-s-1} - \frac{1}{n} \varphi(t) \varphi(\tau),
\end{aligned}
}
where, however, it is to be noted that $t_{n-1} = 0$ is to be assumed, and hence
\[
E_{s,s} = a_1 t_{n-2} + \cdots + a_s t_{n-s-1}
\]
is to be set, while the remaining $E$ are determined by the formulas (9) of the
preceding section.
